\section{Anatomie des trades marquants}

Les deux figures qui suivent isolent les quatre meilleures et les quatre pires
positions. Le trait gris est le cours de clôture quotidien de l'or~; le trait
coloré relie les deux prix réellement obtenus à l'exécution. Les prix du
courtier peuvent s'écarter légèrement du cours de référence, en particulier
lorsqu'une position s'ouvre ou se ferme en cours de séance.

\subsection{Les quatre qui ont fait le résultat}

\begin{figure}[H]
\centering
\includegraphics[width=\linewidth]{gold_trades_best4}
\caption{Les quatre plus gros gains. Ensemble, ils rapportent plus que le
résultat net total du moteur.}
\end{figure}

Le trade du \textbf{21~août~2025}, tenu \metric{214~jours}, rapporte à lui seul
$+$21\,652\,\$ --- soit les deux tiers du résultat net. Le prix d'entrée est
3\,339\,\$ l'once, celui de sortie 4\,405\,\$~: un gain de prix de 31,9\,\%,
sur lequel le portage prélève 2\,874\,\$.

Les trois autres suivent le même schéma~: entrée discrète, tenue de plusieurs
mois, sortie quand le score bascule. Aucun n'a été «~bien~» chronométré au sens
où l'entrée aurait capté un point bas --- sur les quatre panneaux, le cours
monte souvent bien avant l'entrée et continue parfois après la sortie. Le moteur
ne prétend pas attraper les extrêmes~: il prend le milieu du mouvement.

\begin{insightbox}
Sur les quatre panneaux, la sortie intervient nettement après le sommet. C'est
inhérent à un signal de momentum, qui a besoin d'une dégradation confirmée pour
basculer. Ce retard à la sortie est le coût structurel de la méthode, et il est
payé sur chaque grand trade.
\end{insightbox}

\subsection{Les quatre plus coûteux}

\begin{figure}[H]
\centering
\includegraphics[width=\linewidth]{gold_trades_worst4}
\caption{Les quatre plus grosses pertes. La plus lourde est la seule sortie
déclenchée par le stop de sécurité.}
\end{figure}

La perte la plus lourde, $-$1\,955\,\$, est celle du \textbf{1\textsuperscript{er}~avril~2026}~:
la position est ouverte à 4\,759\,\$ et coupée le lendemain à 4\,565\,\$, soit
$-$4,1\,\%. C'est le \textbf{seul} déclenchement du stop de sécurité sur toute
la période, et il se voit dans l'horodatage --- la sortie tombe à 06{:}08~UTC
au lieu de la clôture de séance habituelle.

Ce cas mérite d'être regardé de près, car il illustre ce que le stop fait et ne
fait pas. D'abord, il s'est déclenché sur un creux de séance que le cours de
clôture ne montre pas~: la sortie tombe à 4\,565\,\$ alors que l'or clôture ce
jour-là à 4\,676\,\$. Un examen mené sur données journalières ne verrait tout
simplement pas cet événement.

Ensuite, il a coupé au plus mauvais moment~: le cours revient à 4\,720\,\$ six
jours plus tard et dépasse 4\,840\,\$ à la mi-avril. Le stop n'a donc pas
protégé le résultat sur ce trade précis --- il l'a dégradé. Sa fonction est
ailleurs~: il borne la perte maximale d'une position unique en cas de
décrochage violent, ce qu'aucun autre mécanisme du moteur ne fait. Sur un
moteur qui utilise jusqu'à 6,6$\times$ de levier, c'est une assurance dont on
paie la prime --- ici 1\,955\,\$ pour un seul déclenchement en quatre ans et
demi.

Les trois autres pertes sont ordinaires~: des positions ouvertes sur un score
tout juste positif, refermées en quelques jours quand le mouvement ne se
confirme pas. Aucune ne dépasse 1\,150\,\$. C'est le fonctionnement nominal ---
le moteur teste, se trompe, et sort vite.
